\documentclass[../DoAn.tex]{subfiles}
\begin{document}

\section{Kiến trúc tổng thể}

Repo GEO được tổ chức như một benchmark độc lập cho UAV reconstruction. Hệ thống gồm các thành phần chính:
\begin{enumerate}[leftmargin=1.5em]
  \item module chuẩn bị và đọc dataset;
  \item module predictor cho DUSt3R và MASt3R;
  \item module risk estimation;
  \item module benchmark và visualization;
  \item bộ script bootstrap, POC và full run.
\end{enumerate}

\section{Cấu trúc mã nguồn}

\begin{table}[t]
\centering
\small
\begin{tabularx}{\textwidth}{>{\raggedright\arraybackslash}p{4cm} X}
\toprule
\textbf{Tệp/thư mục} & \textbf{Vai trò} \\
\midrule
\texttt{dataset.py} & Loader cho Dronescapes, ODMData và phần tóm tắt dataset \\
\texttt{realdata.py} & Chuẩn bị dữ liệu thật, export subset, mapping suite dataset \\
\texttt{predictors.py} & Wrapper cho DUSt3R, MASt3R và các output depth \\
\texttt{risk.py} & Tính điểm rủi ro và chọn frame refine \\
\texttt{benchmark.py} & Giao thức benchmark, metric, tổng hợp báo cáo \\
\texttt{scripts/} & Bootstrap, runner POC, runner full thermal-safe \\
\texttt{run\_geo\_project.sh} & Entry point end-to-end cho dự án \\
\bottomrule
\end{tabularx}
\caption{Cấu trúc thành phần chính của hệ thống GEO.}
\label{tab:repo-structure}
\end{table}

\section{Thiết lập thực thi}

\begin{table}[t]
\centering
\small
\begin{tabular}{lp{8.3cm}}
\toprule
Thiết lập & Giá trị \\
\midrule
Coarse image size & 224 \\
Refine image size & 384 \\
Batch size & 1 trong run POC ổn định được dùng cho báo cáo \\
Top-$K$ refined frames & 8 \\
Risk neighbors & 4 \\
Thiết bị coarse/refine & CUDA trên máy server run chính \\
Mã nguồn & \url{https://github.com/duy12i1i7/GEO} \\
\bottomrule
\end{tabular}
\caption{Một số tham số chính của implementation được dùng trong thực nghiệm báo cáo.}
\label{tab:implementation-thesis}
\end{table}

\section{Các chế độ chạy}

Hiện hệ thống hỗ trợ hai kiểu chạy hữu ích:
\begin{enumerate}[leftmargin=1.5em]
  \item \textbf{POC/balanced}: phục vụ kiểm tra ý tưởng trên máy Ubuntu + NVIDIA tầm trung, có thể bỏ OpenMVS để giảm tải;
  \item \textbf{full}: chạy suite rộng hơn cho ODMData và nhiều split hơn của Dronescapes trên máy mạnh.
\end{enumerate}

Thiết kế này giúp đồ án vừa có tính thực dụng trong phát triển, vừa có đường mở để mở rộng benchmark sau này.

\end{document}
